\scanpage{100}
% archive foliation: 99
indexée par $\underline\omega\times\underline t^*$ soit une famille d'éléments
de $V_J(k)$ s'exprime par la formule
\[
(75)\qquad\sum_{i\in J}\psi_i^\omega(\delta)=0
\]
les \deleted{\unsure{2}} conditions (71) et \unsure{(71)|(72)} s'expriment
aussitôt en termes des $\psi_i^\omega(\delta)$, comme
\[
(76)\qquad\psi_i^{\omega'}(\delta)=i_\omega\psi_i^\omega(\delta)
\]
\[
(77)\qquad\psi_i^\omega(\lambda\delta)=\lambda^{-1}\psi_i^\omega(\delta).
\]
Il reste à exprimer le fait que les éléments
$\psi^\omega(\delta)\in V_J(k)$ définis par les $(\omega,\delta)\in\underline\omega\times\underline t_{\unsure{E}}^*$ sont
dans $F_t$ --- moyennant (76) (77) il suffit de
l'exprimer pour un seul couple $(\omega,\delta)$ ---,
et la condition (73).
\deleted{\unsure{En exprimant ces deux} \illegible{1} \unsure{simultanément, en écrivant}}
\inserted{Si on définit $F_t$ comme noyau d'une forme linéaire sur $V_J$}
\[
\deleted{(77)\quad\sum_{i\in J}\psi_i^\omega(\delta)e_i\unsure{-|=}}
\]
\[
F_t=\deleted{\unsure{\Sigma}}\left\{\sum\lambda_i e_i\ \middle|\ \sum\lambda_i=0,\quad
\sum c_i\lambda_i\deleted{\unsure{e_i}}=0\right\}
\]
% margin position: left
\margin{$(e_i)_{i\in J}$ base canonique de $k^J$}
[\uline{NB} l'élément $(c_i)_{i\in J}$ de \deleted{$\unsure{\Sigma}$}$k^J$ est déterminé
modulo multiplication par $k^*$ et addition d'un élément
diagonal], \deleted{cette} condition $\psi_i^\omega(\delta)\in F_t$
s'exprime par la condition calculatoire
\[
(78)\qquad\sum_{i\in J}c_i\psi_i^\omega(\delta)=0.
\]
Il reste \inserted{ensuite} à exprimer (73)
